%%% BAC solution 06 %%%
\documentclass{article}
\usepackage[margin=1in]{geometry}
\usepackage{amsmath, amssymb, amsfonts, enumitem, fancyhdr, tikz}
% get rid of paragraph indent
\setlength{\parindent}{0 pt}
% allow section.equation numbering
\numberwithin{equation}{section}
% allows you to copy-paste code segments. requires pygments, which can be
% installed from PyPI with pip install Pygments.
% warning: minted does NOT work with Texmaker if you have TeX Live installed
% on WSL but Texmaker installed natively on Windows!
%\usepackage{minted}
% alternative to minted that does not require Python, LaTeX only. listings is
% however disgusting out of the box and some setup is required.
\usepackage{listings, xcolor}
% makes clickable links to sections
\usepackage{hyperref}
% make the link colors blue, as well as cite colors. urls are magenta
\hypersetup{
    colorlinks, linkcolor = blue, citecolor = blue, urlcolor = magenta
}
% fancy pagestyle so we can use fancyhdr for fancy headers/footers
\pagestyle{fancy}
% add logo in right of header. note that you will have to adjust logo path!
\fancyhead[R]{\includegraphics[scale = 0.15]{../bac_logo1.png}}
% don't show anything in the left and center header
\fancyhead[L, C]{}
% give enough space for logo by reducing top margin height, head separator,
% increasing headerheight. see Figure 1 in the fancyhdr documentation. if
% \topmargin + \headheight + \headsep = 0, original text margins unchanged.
\setlength{\topmargin}{-60 pt}
\setlength{\headheight}{50 pt}
\setlength{\headsep}{10 pt}
% remove decorative line in the fancy header
\renewcommand{\headrulewidth}{0 pt}

% color definitions for listings syntax highlighting. uses colors borrowed
% from the VS Code Dark+ and Abyss standard themes.
\definecolor{KwColor}{RGB}{153, 102, 184}     % keyword color
\definecolor{VarColor}{RGB}{86, 156, 214}     % variables/identifier color
\definecolor{StrColor}{RGB}{209, 105, 105}    % string color
\definecolor{CmtColor}{RGB}{106, 153, 85}     % comment color

% general listings configuration for all languages
\lstset{
    % change keyword, identifier, comment, string colors
    keywordstyle = \color{KwColor},
    commentstyle = \color{CmtColor},
    identifierstyle = \color{VarColor},
    stringstyle = \color{StrColor},
    % no spaces in strings
    showstringspaces = false,
    % monospace font by default
    basicstyle = \ttfamily,
    % tabsize 8 by default, this is not the 1960s
    tabsize = 4,
    % add line numbers to the left with gray typewriter font
    numbers = left,
    numberstyle = \color{gray}\ttfamily,
    % change distance from code block from 10 pt to 5 pt
    numbersep = 5 pt
}

% title, author + thanks, date
\title{Solution 6}
\author{Derek Huang\thanks{NYU Stern 2021, BAC Advanced Team.}}
% the extra thanks is also useful for making notes about previous versions.
% you can remove this thanks if there are no version revisions.
\date{April 19, 2021}

\begin{document}

\maketitle
% need to include this after making title to undo the automatic
% \thispagestyle{plain} command that is issued.
\thispagestyle{fancy}

% remove some space between title and first section to keep text on one page.
% comment this line out or adjust the space reduction as needed. generally,
% try to remove as little space as possible if you are trying to keep content
% on one page, else comment this line out if content spans multiple pages.
%\vspace{-20 pt}

\section{Convexity}

\subsection{The convex hull}

\textit{Proof.} Suppose $ C $ is convex. Then, $ \forall \mathbf{x}_1,
\mathbf{x}_2 \in C $, $ \alpha_1 \in [0, 1] $, $ \gamma_{1, 2} \triangleq
\alpha_1\mathbf{x}_1 + (1 - \alpha_1)\mathbf{x}_2 \in C $. Choose an arbitrary
$ \mathbf{x}_3 \in C $, $ \alpha_2 \in [0, 1] $. Then, $ C $ convex,
$ \gamma_{1, 2} \in C \Rightarrow \gamma_{1, 3} \triangleq
\alpha_2\mathbf{x}_3 + (1 - \alpha_2)\gamma_{1, 2} \in C $. $ \forall n \in
\mathbb{N} $, $ \alpha_{n - 1} \in [0, 1] $, $ \mathbf{x}_n \in C $,
$ \gamma_{1, n} \triangleq \alpha_{n - 1}\mathbf{x}_n + (1 - \alpha_{n - 1})
\gamma_{1, n - 1} \in C $, which can be written as the convex combination
\begin{equation*}
    \alpha_{n - 1}\mathbf{x}_n + \left(
        \prod_{j = 2}^{n - 1}(1 - \alpha_j)
    \right)\left(
        \alpha_1\mathbf{x}_1 + (1 - \alpha_1)\mathbf{x}_2 +
        \sum_{k = 3}^{n - 1}\alpha_{k - 1}\mathbf{x}_k
    \right)
\end{equation*}

Since $ \mathbf{x}_1, \ldots \mathbf{x}_n \in C $, $ \alpha_1, \ldots
\alpha_{n - 1} \in [0, 1] $ arbitrary, then any $ \gamma_{1, n} \in C $ is a
convex combination, so $ C $ contains all convex combinations of its elements,
i.e. $ \operatorname{conv} C \subseteq C $. Since $ C \subseteq
\operatorname{conv} C $, we conclude that $ C = \operatorname{conv} C $. Now,
suppose that $ C = \operatorname{conv} C $. Then, $ \forall \mathbf{x}_1,
\mathbf{x}_2 \in C $, $ \alpha \in [0, 1] $, $ \gamma_{1, 2} \triangleq
\alpha\mathbf{x}_1 + (1 - \alpha)\mathbf{x}_2 \in \operatorname{conv} C $, and
since $ C = \operatorname{conv} C $, $ \gamma_{1, 2} \in C \Rightarrow C $
convex, by definition.

\subsection{Positive definite matrices}

\textit{Proof.} Choose arbitrary $ \mathbf{A}, \mathbf{B} \in
\mathbb{S}^n_{\succ \mathbf{0}} $. $ \mathbb{S}^n_{\succ \mathbf{0}} $ closed
under multiplication by positive scalars $ \Rightarrow \forall \alpha
\in (0, 1) $, we have $ \alpha\mathbf{A}, (1 - \alpha)\mathbf{B} \in
\mathbb{S}^n_{\succ \mathbf{0}} $. $ \mathbb{S}^n_{\succ \mathbf{0}} $ closed
under matrix addition $ \Rightarrow \forall \alpha \in [0, 1] $, we have
$ \alpha\mathbf{A} + (1 - \alpha)\mathbf{B} \in
\mathbb{S}^n_{\succ \mathbf{0}} $. $ \mathbf{A}, \mathbf{B}, $ arbitrary
$ \Rightarrow \forall \alpha \in [0, 1], \forall \mathbf{A}, \mathbf{B} \in
\mathbb{S}^n_{\succ \mathbf{0}} $, $ \alpha\mathbf{A} +
(1 - \alpha)\mathbf{B}) \in \mathbb{S}^n_{\succ \mathbf{0}} \Rightarrow
\mathbb{S}^n_{\succ \mathbf{0}} $ convex.

\subsection{Norms}

\textit{Proof.} Choose arbitrary $ \mathbf{u}, \mathbf{v} \in V $, $ \alpha \in
[0, 1] $. By the triangle inequality and absolute homogeneity properties of
norms, $ \Vert\alpha\mathbf{u} + (1 - \alpha)\mathbf{v}\Vert \le
|\alpha|\Vert\mathbf{u}\Vert + |1 - \alpha|\Vert\mathbf{v}\Vert =
\alpha\Vert\mathbf{u}\Vert + (1 - \alpha)\Vert\mathbf{v}\Vert $. Since
$ \alpha \in [0, 1] $, $ \mathbf{u}, \mathbf{v} \in V $ arbitrary, the
inequality holds $ \forall \alpha \in [0, 1], \forall \mathbf{u}, \mathbf{v}
\in V \Rightarrow V $ is convex.

% no \hrulefill since the footnote line is there

\subsection{Scalar composition}

\textit{Proof.} $ f $ nondecreasing $ \Rightarrow \forall x, y \in
\mathbb{R} $, $ x < y $, that $ f(x) \le f(y) $. $ g $ convex $ \Rightarrow
\forall \mathbf{x}, \mathbf{y} \in \mathbb{R}^n $, $ \alpha \in [0, 1] $,
that $ g(\alpha\mathbf{x} + (1 - \alpha)\mathbf{y}) \le
\alpha g(\mathbf{x}) + (1 - \alpha)g(\mathbf{y}) $. Since $ f $ is also
convex, then
\begin{equation*}
    f \circ g(\alpha\mathbf{x} + (1 - \alpha)\mathbf{y}) \le
    f(\alpha g(\mathbf{x}) + (1 - \alpha)g(\mathbf{y})) \le
    \alpha f \circ g(\mathbf{x}) + (1 - \alpha)f \circ g(\mathbf{y})
\end{equation*}

This holds $ \forall \alpha \in [0, 1] $, $ \forall \mathbf{x}, \mathbf{y} \in
\mathbb{R}^n $, so $ f \circ g $ is convex.

\section{Optimization}

\subsection{Underdetermined linear systems}

\textit{Solution.} The Lagrangian $ \mathcal{L}_f : \mathbb{R}^n \times
\mathbb{R}^q \rightarrow \mathbb{R} $ is such that
\begin{equation*}
    \mathcal{L}_f(\mathbf{x}, \nu) \triangleq \mathbf{x}^\top\mathbf{x} +
    \nu^\top(\mathbf{Ax} - \mathbf{b})
\end{equation*}

Strong duality holds, so $ \exists \mathbf{x}^* \in \mathbb{R}^n $ primal
optimal, $ \exists \nu^* \in \mathbb{R}^q $ dual optimal. Letting $ f_d $
denote the dual of the objective, strong duality $ \Rightarrow
\Vert\mathbf{x}^*\Vert_2^2 = f_d(\nu^*) \triangleq \inf_{\mathbf{x} \in
\mathbb{R}^n}\{\mathcal{L}_f(\mathbf{x}, \nu^*)\} =
\mathcal{L}_f(\mathbf{x}^*, \nu^*) $. Since $ \mathcal{L}_f $ convex in
$ \mathbf{x} $,
\begin{equation*}
    \nabla_\mathbf{x}\mathcal{L}_f(\mathbf{x}^*, \nu^*) =
    2\mathbf{x}^* + \mathbf{A}^\top\nu^* = \mathbf{0} \Leftrightarrow
    \mathbf{x}^* = -\frac{1}{2}\mathbf{A}^\top\nu^*
\end{equation*}

Therefore, substituting our expression for $ \mathbf{x}^* $ into the
definition of $ \mathcal{L}_f $, we can write the dual $ f_d $ at $ \nu^* $ as
\begin{equation*}
    f_d(\nu^*) = -\frac{1}{4}\big\Vert\mathbf{A}^\top\nu^*\big\Vert_2^2 -
    \mathbf{b}^\top\nu^*
\end{equation*}

Since $ f_d $ concave, then $ \nabla f_d(\nu^*) = \mathbf{0} $, and we see that
\begin{equation*}
    \nabla f_d(\nu^*) = -\frac{1}{2}\mathbf{AA}^\top\nu^* - \mathbf{b} =
    \mathbf{0} \Leftrightarrow
    \nu^* = -2\big(\mathbf{AA}^\top\big)^{-1}\mathbf{b}
\end{equation*}

Substituting this expression into the previous expression for $ \mathbf{x}^* $,
we see that
\begin{equation*}
    \mathbf{x}^* = \mathbf{A}^\top\big(\mathbf{AA}^\top\big)^{-1}\mathbf{b}
\end{equation*}

\subsection{Markowitz portfolios}

\textit{Solution.} The Lagrangian $ \mathcal{L}_f : \mathbb{R}^n \times \mathbb{R} \rightarrow \mathbb{R} $ is such that
\begin{equation*}
    \mathcal{L}_f(\mathbf{w}, \nu) \triangleq \mathbf{w}^\top\mathbf{Cw} +
    \nu\big(\mathbf{1}^\top\mathbf{w} - 1\big)
\end{equation*}

Strong duality holds, so $ \exists \mathbf{w}^* \in \mathbb{R}^n $ primal
optimal, $ \exists \nu^* \in \mathbb{R} $ dual optimal. Strong duality
$ \Rightarrow p^* = f_d(\nu^*) \triangleq \inf_{\mathbf{w} \in \mathbb{R}^n}
\{\mathcal{L}(\mathbf{w}, \nu^*)\} = \mathcal{L}(\mathbf{w}^*, \nu^*) $.
Since $ \mathcal{L}_f $ is convex in $ \mathbf{w} $, then
\begin{equation*}
    \nabla\mathcal{L}_f(\mathbf{w}^*, \nu^*) = 2\mathbf{Cw}^* +
    \nu^*\mathbf{1} = \mathbf{0} \Leftrightarrow
    \mathbf{w}^* = -\frac{1}{2}\nu^*\mathbf{C}^{-1}\mathbf{1}
\end{equation*}

Since $ \nu^* $ is scalar, we can use the equality constraint to see that
\begin{equation*}
    1 = \mathbf{1}^\top\mathbf{w}^* =
    -\frac{1}{2}\nu^*\mathbf{1}^\top\mathbf{C}^{-1}\mathbf{1} \Leftrightarrow
    \nu^* = -\frac{2}{\mathbf{1}^\top\mathbf{C}^{-1}\mathbf{1}}
\end{equation*}

Substituting the result for $ \nu^* $ into the previous expression for
$ \mathbf{w}^* $, we have
\begin{equation*}
    \mathbf{w}^* = \frac{
        \mathbf{C}^{-1}\mathbf{1}
    }{
        \mathbf{1}^\top\mathbf{C}^{-1}\mathbf{1}
    }
\end{equation*}

\subsection{Highly constrained portfolios}

\textit{Solution.} It is instructive to first write the problem in standard
form, i.e. as
\begin{equation*}
    \begin{array}{ll}
        \displaystyle\min_\mathbf{w} & \mathbf{w}^\top\mathbf{Cw} \\
        \text{s.t.} & -\mathbf{w} \preceq \mathbf{0} \\
            & \mathbf{w} - \gamma\mathbf{1} \preceq \mathbf{0} \\
            & -\mathbf{1}^\top\mathbf{w} \le 0 \\
            & \mathbf{1}^\top\mathbf{w} - (1 - \kappa) \le 0
    \end{array}
\end{equation*}

Here we choose to write the scalar inequalities last for ease of
identification. In total, we can see that there are $ 2n + 2 $ inequality
constraints, so the Lagrangian $ \mathcal{L}_f : \mathbb{R}^n \times
\mathbb{R}^{2n + 2} \rightarrow \mathbb{R} $ can be written such that
\begin{equation*}
    \mathcal{L}_f(\mathbf{w}, \lambda) \triangleq \mathbf{w}^\top\mathbf{Cw} +
        \lambda^\top\begin{bmatrix}
            \ -\mathbf{w} \ \\
            \ \mathbf{w} - \gamma\mathbf{1} \\ \
            \ -\mathbf{1}^\top\mathbf{w} \ \\
            \ \mathbf{1}^\top\mathbf{w} - (1 - \kappa) \
        \end{bmatrix}
\end{equation*}

Since $ \mathcal{L}_f $ is convex in $ \mathbf{w} $, $ \forall \lambda \in
\mathbb{R}^{2n + 2} $, at $ \mathbf{w}_\lambda \triangleq
\arg\min_{\mathbf{w} \in\mathbb{R}^n}\mathcal{L}_f(\mathbf{w}, \lambda) $,
$ \nabla_\mathbf{w}\mathcal{L}_f(\mathbf{w}_\lambda, \lambda) = \mathbf{0} $. Expanding,
\begin{equation} \label{p2.3_wlambda}
    \nabla_\mathbf{w}\mathcal{L}_f(\mathbf{w}_\lambda, \lambda) =
    2\mathbf{Cw}_\lambda + \begin{bmatrix}
        \ -\mathbf{I} & \mathbf{I} & -\mathbf{1} & \mathbf{1} \
    \end{bmatrix}\lambda = \mathbf{0} \Leftrightarrow
    \mathbf{w}_\lambda = -\frac{1}{2}\begin{bmatrix}
        \ -\mathbf{C}^{-1} & \mathbf{C}^{-1} & -\mathbf{C}^{-1}\mathbf{1} &
        \mathbf{C}^{-1}\mathbf{1} \
    \end{bmatrix}\lambda
\end{equation}

Here the matrix $ [ \ -\mathbf{I} \ \ \mathbf{I} \ \ -\mathbf{1} \ \
\mathbf{1} \ ] \in \mathbb{R}^{n \times (2n + 2)} $, as $ \mathbf{I} \in
\mathbb{R}^{n \times n} $, $ \mathbf{1} \in \mathbb{R}^n $. The above result
can be more clearly seen if the inequality constraints of $ \mathcal{L}_f $
are written more explicitly. We can write
\begin{equation*}
    \begin{split}
    \mathcal{L}_f(\mathbf{w}, \lambda) & = \mathbf{w}^\top\mathbf{Cw} +
    \sum_{k = 1}^n(
        -\lambda_kw_k + \lambda_{n + k}(w_k - \gamma) - \lambda_{2n + 1}w_k +
        \lambda_{2n + 2}w_k
    ) - \lambda_{2n + 2}(1 - \kappa) \\ & =
    \mathbf{w}^\top\mathbf{Cw} + \sum_{k = 1}^n\underbrace{
        (-\lambda_k + \lambda_{n + k} - \lambda_{2n + 1} + \lambda_{2n + 2})
    }_{
        k\text{th row of } [
            \ -\mathbf{I} \ \ \mathbf{I} \ \ -\mathbf{1} \ \ \mathbf{1} \
        ]\lambda
    }w_k - \gamma\sum_{k' = 1}^n\lambda_{n + k'} - \lambda_{2n + 2}(1 - \kappa)
    \end{split}
\end{equation*}

By the definition of the dual, $ f_d(\lambda) \triangleq \inf_{\mathbf{w}
\in \mathbb{R}^n}\{\mathcal{L}_f(\mathbf{w}, \lambda)\} =
\mathcal{L}_f(\mathbf{w}_\lambda, \lambda) $, so substituting, we have
\begin{equation*}
    \begin{split}
    f_d(\lambda) = \frac{1}{4}\lambda^\top\begin{bmatrix}
        -\mathbf{I} \ \\ \phantom{-} \mathbf{I} \ \\ \ \ -\mathbf{1}^\top \ \\
        \ \ \phantom{-} \mathbf{1}^\top \
    \end{bmatrix}\mathbf{C}^{-1}\begin{bmatrix}
        \ -\mathbf{I} & \mathbf{I} & -\mathbf{1} & \mathbf{1} \
    \end{bmatrix}\lambda & - \frac{1}{2}\lambda^\top\begin{bmatrix}
        % row 1
        \ \phantom{-}\mathbf{C}^{-1} & -\mathbf{C}^{-1} &
        \phantom{-}\mathbf{C}^{-1}\mathbf{1} & -\mathbf{C}^{-1}\mathbf{1} \ \\
        % row 2
        \ -\mathbf{C}^{-1} & \phantom{-}\mathbf{C}^{-1} & 
        -\mathbf{C}^{-1}\mathbf{1} & \phantom{-}\mathbf{C}^{-1}\mathbf{1} \ \\
        % row 3
        \phantom{-}\mathbf{1}^\top\mathbf{C}^{-1} &
        -\mathbf{1}^\top\mathbf{C}^{-1} &
        \phantom{-}\mathbf{1}^\top\mathbf{C}^{-1}\mathbf{1} &
        -\mathbf{1}^\top\mathbf{C}^{-1}\mathbf{1} \ \\
        % row 4
        -\mathbf{1}^\top\mathbf{C}^{-1} &
        \phantom{-}\mathbf{1}^\top\mathbf{C}^{-1} &
        -\mathbf{1}^\top\mathbf{C}^{-1}\mathbf{1} &
        \phantom{-}\mathbf{1}^\top\mathbf{C}^{-1}\mathbf{1} \
    \end{bmatrix}\lambda \\
    % starts at the subtraction of second block matrix term
    & + \lambda^\top\begin{bmatrix}
        \ \mathbf{0} \ \\ -\gamma\mathbf{1} \ \\ \ 0 \ \\ -(1 - \kappa) \
    \end{bmatrix}
    \end{split}
\end{equation*}

Multiplying the block matrices out in the first term, we can make the note that
\begin{equation*}
    \begin{bmatrix}
        -\mathbf{I} \ \\ \phantom{-} \mathbf{I} \ \\ \ \ -\mathbf{1}^\top \ \\
        \ \ \phantom{-} \mathbf{1}^\top \
    \end{bmatrix}\mathbf{C}^{-1}\begin{bmatrix}
        \ -\mathbf{I} & \mathbf{I} & -\mathbf{1} & \mathbf{1} \
    \end{bmatrix} = \begin{bmatrix}
        % row 1
        \ \phantom{-}\mathbf{C}^{-1} & -\mathbf{C}^{-1} &
        \phantom{-}\mathbf{C}^{-1}\mathbf{1} & -\mathbf{C}^{-1}\mathbf{1} \ \\
        % row 2
        \ -\mathbf{C}^{-1} & \phantom{-}\mathbf{C}^{-1} & 
        -\mathbf{C}^{-1}\mathbf{1} & \phantom{-}\mathbf{C}^{-1}\mathbf{1} \ \\
        % row 3
        \phantom{-}\mathbf{1}^\top\mathbf{C}^{-1} &
        -\mathbf{1}^\top\mathbf{C}^{-1} &
        \phantom{-}\mathbf{1}^\top\mathbf{C}^{-1}\mathbf{1} &
        -\mathbf{1}^\top\mathbf{C}^{-1}\mathbf{1} \ \\
        % row 4
        -\mathbf{1}^\top\mathbf{C}^{-1} &
        \phantom{-}\mathbf{1}^\top\mathbf{C}^{-1} &
        -\mathbf{1}^\top\mathbf{C}^{-1}\mathbf{1} &
        \phantom{-}\mathbf{1}^\top\mathbf{C}^{-1}\mathbf{1} \
    \end{bmatrix}
\end{equation*}

We can therefore combine the first two terms in the expression for
$ f_d(\lambda) $ and arrive at
\begin{equation*}
    f_d(\lambda) = -\frac{1}{4}\lambda^\top\begin{bmatrix}
        -\mathbf{I} \ \\ \phantom{-} \mathbf{I} \ \\ \ \ -\mathbf{1}^\top \ \\
        \ \ \phantom{-} \mathbf{1}^\top \
    \end{bmatrix}\mathbf{C}^{-1}\begin{bmatrix}
        \ -\mathbf{I} & \mathbf{I} & -\mathbf{1} & \mathbf{1} \
    \end{bmatrix}\lambda + \lambda^\top\begin{bmatrix}
        \ \mathbf{0} \ \\ -\gamma\mathbf{1} \ \\ \ 0 \ \\ -(1 - \kappa) \
    \end{bmatrix}
\end{equation*}

By properties of the dual, $ f_d $ is concave, so the block matrix product
must be positive semidefinite, as $ f_d $ is a quadratic function. Furthermore,
$ \forall \lambda \in \mathbb{R}^{2n + 2} $, $ f_d(\lambda) \in \mathbb{R} $,
which means $ \operatorname{dom} f_d = \mathbb{R}^{2n + 2} $. Therefore, the
dual problem for the primal portfolio optimization problem is simply
\begin{equation*}
    \begin{array}{ll}
        \displaystyle\max_\lambda & f_d(\lambda) \\
        \text{s.t.} & \lambda \succeq \mathbf{0}
    \end{array}
\end{equation*}

Since the primal problem is a convex optimization problem with affine
constraints and open domain\footnote{
    $ \mathbb{R}^n $ is both open and closed.
}, assuming that the feasible set is nonempty, then $ \exists \mathbf{w}^* \in
\mathbb{R}^n $ primal optimal, $ \exists \lambda^* \succeq \mathbf{0} \in
\mathbb{R}^{2n + 2} $ dual optimal. The $ \lambda^* $ obtained from solving
the dual problem can then be used in (\ref{p2.3_wlambda}) to compute
$ \mathbf{w}^* $.

\end{document}